% The five-wave planar Riemann solver.  Written 2026-09-08 from the code and
% the measurements as they stand (HLLD @ 954046b, rmhd_final @ add9f70).
% Every number here is measured; docs/planar_solver_notes.md gives the file,
% line, or script behind each one.
\section{A Planar Riemann Solver}
\label{sec:planar}

The seven-wave solver of Sect.~\ref{sec:setup:exact} resolves only a minority
of the interfaces a two-dimensional run presents. This section explains why,
and describes the reduced solver that follows from the explanation. The
reduction is not an approximation: a converged answer is verified against the
complete seven-wave jump conditions before it is accepted.

\subsection{Why the seven-wave form fails on a planar flow}
\label{sec:planar:why}

A two-dimensional flow keeps the magnetic field and the velocity in a single
plane. On a magnetic rotor this is not approximate: measured over all
interfaces of a $64^2$ run, the coplanarity of the two tangential fields,
%
\begin{equation}
  \chi \;=\;
  \frac{\lvert \mathbf{B}_t^L \times \mathbf{B}_t^R \rvert}
       {\lvert \mathbf{B}_t^L \rvert \, \lvert \mathbf{B}_t^R \rvert},
  \label{eq:coplanarity}
\end{equation}
%
has median $3.5 \times 10^{-11}$ and lies below $10^{-6}$ on every interface.
For comparison, a forward-constructed dataset drawn from the same envelope has
0.01\% of samples below $10^{-4}$ and a median of $0.65$: the geometry of a two-dimensional run is not a
sampling corner, it is a different population.

In that geometry the seven-wave parameterisation is ill-posed in three
separate ways, each measured on the $481\,837$ interfaces the run solved
exactly.

\paragraph{An unknown that no wave determines.}
A planar slow wave cannot rotate the tangential field; it can only reverse it.
Measured, the left slow wave turns the field by a median of $6 \times 10^{-12}$
radians, which is zero to eleven digits. The orientation $\psi^\ast$ is
therefore fixed by the geometry, while the solver still treats it as an
unknown and differences the residual with respect to it. That column of the
Jacobian is finite-difference noise. Freezing the three angles at their true
values raises the converged fraction from $79.5\%$ to $90.5\%$, which is the
size of the damage.

\paragraph{A wave asked for two numbers that can supply one.}
Each slow wave is parameterised by the tangential field it must produce at the
contact, which is two components, while the wave family has one parameter.
This is why the formulation carries a slack residual per slow wave. The
over-determination would be harmless if the waves were strong, but they are
not: the left slow wave changes $\lvert\mathbf{B}_t\rvert$ by a median of only
$3.3\%$, so the set of contact fields it can reach is small and most trial
values of the unknowns lie outside it. The wave construction then fails
outright rather than returning a large residual. This is visible in the
failure pattern: only about half of the failed interfaces can be constructed
at all from a given starting point, while $99.5\%$ can be constructed
somewhere inside a wide grid. The search is landing outside a feasible
region, not climbing a hill.

\paragraph{Two nearly parallel wave families.}
The Alfv\'en and slow speeds differ by a median of $5.5 \times 10^{-3}$, with
a fifth percentile of $9.4 \times 10^{-6}$, against $2.8 \times 10^{-1}$
separating the fast wave from the Alfv\'en wave. The seven-wave system therefore
carries two almost coincident characteristic fields wherever the rotor is
concerned.

The consequence is that no starting point rescues these interfaces. Since a
planar flow admits only eight discrete choices for the three angles, the
orientation being the common line or its reverse and each rotation being
nothing or a reversal, that space can be searched exhaustively. Crossing all
eight with a $7^3$ grid over the three remaining unknowns, $2744$ points per
interface, converges on $8.5\%$; a grid a hundred times denser than $3^3$
gains $1.5$ percentage points and does not lower the best residual at all.

\subsection{The reduced system}
\label{sec:planar:system}

In the plane the state is $(\rho,\,p_{\rm tot},\,v^x,\,v^t,\,B^t)$ with $B^n$
fixed and $B^t$ a \emph{signed} scalar. The system has five characteristic
fields, so the Riemann fan has five waves: two fast, two slow, and the
contact. There are no rotational discontinuities, and therefore no angles and
no nearly parallel pair. Reading forward from each side,
%
\begin{equation}
  \mathbf{u} = \bigl[\, \ln p^\ast_{\mathrm{LF}},\;
                        B^{t\ast},\;
                        \ln p^\ast_{\mathrm{RF}} \,\bigr],
  \qquad
  \mathbf{f} = \bigl[\, \llbracket v^x \rrbracket,\;
                        \llbracket v^t \rrbracket,\;
                        \llbracket p_{\rm tot} \rrbracket \,\bigr],
  \label{eq:planar-system}
\end{equation}
%
a square $3 \times 3$ system. Each slow wave is given exactly one number to
reach, which is its family's dimension, so the slack equations disappear;
they are retained only as a diagnostic, and come out at $10^{-11}$.

Two properties of this parameterisation matter in practice. Because
$B^{t\ast}$ is signed and enters linearly, a reversal of the tangential field
is an ordinary sign change through zero rather than a discrete branch the
solver must be told about, so the eight-way enumeration of angle choices
disappears. And because the system is well conditioned, it needs no warm
start: the results below are obtained from the trivial guess of the mean
pressure on the two sides and the mean tangential field, with no neural
network and no retries.

The solve is performed in the frame rotated about the normal so that the left
tangential field lies along $+\hat{y}$. The out-of-plane residue that this
rotation leaves behind is reported, so an interface that was not planar
enough can be refused; on real interfaces that residue has median $1.5 \times 10^{-10}$ and $0.2\%$ exceed the
tolerance.

\subsection{Verification, and why it is required}
\label{sec:planar:verify}

A five-wave solution is a seven-wave solution whose two rotational
discontinuities have zero strength. It can therefore be written in the
seven-wave unknowns, with both rotations zero and the contact orientation
equal to $0$ or $\pi$ according to the sign of $B^{t\ast}$, and the complete
six-component residual evaluated there. The solver performs this check on
every answer and accepts only those below $10^{-8}$; the median achieved on
real interfaces is $1.1 \times 10^{-10}$.

The check is not a formality. The five-wave family assumes both rotational
discontinuities have zero strength, and a coplanar flow almost enforces that:
a rotation in the plane can only be nothing or a reversal, since those are the
only two rotations carrying the common line to itself. Measured over
$142\,554$ rotations (both sides of $71\,277$ solved interfaces), $97.3\%$
lie below $10^{-6}$ and $0.375\%$ within $10^{-3}$ of $\pi$.

The remaining $2.3\%$ are neither, and they are the informative ones. They sit
on interfaces that are measurably less coplanar: $\chi$ has median
$4.3\times10^{-9}$ there against $7.6\times10^{-11}$ over all interfaces, a
factor of $57$. A departure from coplanarity of one part in $10^{9}$ is
therefore enough to unpin the angle, and $0.27\%$ of rotations exceed $0.1$
radians --- an $O(1)$ response to an $O(10^{-9})$ perturbation. The
restriction to $\{0,\pi\}$ is exact only in the exact coplanar limit; at
finite precision it is a strong tendency with a measurable tail. Where the
true solution does carry a rotation, of $\pi$ or otherwise, the five-wave
family does not contain it and the solver may still converge on its own
reduced system to a different state. The verification is what separates the
cases.

The same test distinguishes this reduction from a cruder one. A three-wave
solver, retaining only the two fast waves and the contact, converges on
essentially every interface, including all of those the seven-wave solver
fails, but not one of its answers reaches $10^{-6}$ on the full residual; the
median is $1.2 \times 10^{-2}$. It is a reliable approximation, not a solution,
and it is not used.

One caveat applies to any such check. The seven-wave residual is
\emph{frame-sensitive} at this tolerance. The orientation of the solver frame
is set arbitrarily by the sweep direction, and the formulation carries the
tangential field as a magnitude and an angle, so rotating a solution changes
how well it appears to satisfy its own equations. Measured on solutions
produced by the planar solver, the fraction satisfying $10^{-8}$ falls from
$100\%$ in the rotated frame to $51\%$ in the raw frame, a tenth of them
failing to construct at all. The effect is present for solutions produced by
the seven-wave solver too, at tens of percent, though its direction there was
not consistent between two harvests and we do not claim one. The safe
practice, and what the solver does internally, is to verify in the frame the
answer was produced in.

\subsection{Uniqueness, and what the tolerance pins}
\label{sec:planar:unique}

A vanishing residual establishes that the answer is \emph{an} admissible
elementary-wave solution, not that it is the only one, and magnetohydrodynamic
Riemann problems are not guaranteed a unique one. We therefore compared full
answers --- every intermediate state and every wave speed, not the star
pressure alone --- from independent routes.

The two seeds of the planar solver have largely disjoint basins: the
$\sqrt{2\bar{p}}$ seed resolves lanes the default cannot and loses $1749$ of
$3619$ that it can, so where both converge they are independent. On $2525$
such interfaces the two answers agree to a median of $7.7\times10^{-12}$ and a
maximum of $9.3\times10^{-10}$, with \emph{no} interface above $10^{-8}$.

Comparing the seven-wave solver with the planar one is more interesting. On
$5095$ interfaces whose harvested solution carries no rotation, the two agree
to a median of $4.3\times10^{-10}$, but $2.0\%$ differ by more than $10^{-6}$
while \emph{both} satisfy the complete jump conditions to $10^{-8}$. This is
not non-uniqueness but conditioning: on precisely those interfaces the
seven-wave Jacobian has condition number $7.9\times10^{3}$ against $32$ where
the two agree, and the product of that condition number with the achieved
residual predicts a state discrepancy of $1.8\times10^{-5}$ against
$1.1\times10^{-5}$ observed. The planar Jacobian's condition number is
$\approx 10$ on both groups and predicts $1.3\times10^{-7}$, two orders too
small. The two solvers find the same solution; the seven-wave
parameterisation localises it some $250$ times less well.

The practical consequence is that a residual tolerance is not a state
tolerance, and the factor between them is the condition number of whichever
parameterisation produced the answer. In the reduced form $10^{-8}$ in the
residual pins the state to $\sim\!10^{-7}$; in the seven-wave form, on the
coplanar interfaces a two-dimensional run supplies, it pins it only to
$\sim\!10^{-4}$ on a small fraction. This is a further reason to prefer the
reduced formulation where it applies: it is not only better conditioned for
\emph{finding} the answer, but for \emph{certifying} it.

\subsection{Measured behaviour}
\label{sec:planar:results}

The planar solver is used as a fallback: an interface is offered to it only
when the seven-wave solver has not already resolved it. Table~\ref{tab:planar}
collects the measurements, all on interfaces harvested from a $64^2$ rotor
run.

\begin{table}
\centering
\caption{The three solvers on the same interfaces. ``Full residual'' is the
complete seven-wave residual evaluated at the solver's own answer.}
\label{tab:planar}
\begin{tabular}{lccc}
\hline
 & seven-wave & three-wave & five-wave planar \\
\hline
converged, interfaces the run resolves & 79.5\% & 99.9\% & 91.2\% \\
converged, interfaces the run fails    & 0.25\% & 99.8\% & 65.8\% \\
full residual below $10^{-8}$          & ---    & 0.0\%  & 96--99.5\% \\
contact pressure vs.\ known solution   & exact  & $10^{-4}$--$4\times10^{-3}$ & $1.6\times10^{-10}$ \\
cost per interface                     & $\sim$30\,ms & $\sim$1\,ms & $\sim$8\,ms \\
warm start required                    & yes    & no     & no \\
\hline
\end{tabular}
\end{table}

Enabling the fallback is strictly additive. On eighteen recorded sweeps
replayed with identical inputs, the number of interfaces receiving an exact
flux rises from $471$ to $763$: $292$ gained, none lost, and on the $471$
resolved by both configurations the star pressure is bitwise identical. In a
$64^2$ run carried to $t = 0.4$ the fraction of attempted interfaces resolved
exactly rises from $43.9\%$ to $78.0\%$, and the improvement is concentrated
where the geometry predicts: the sweeps along which the normal field is
small, previously the weakest at $24.7\%$, reach $65.7\%$. The cost is
roughly a factor of two in wall time per step, since a second solver runs on
every interface the first one leaves.

These are run-level figures, and it is worth stating what they are fractions
of. The weak-jump gate admits $14.1\%$ of interfaces; the exact flux
therefore reaches $11.0\%$ of all interfaces in the run, and HLLD supplies
the rest. Coverage percentages quoted here and elsewhere are of
\emph{attempted} interfaces unless said otherwise.
